
\subsection{Transformaciones de base: cilíndricas, cartesianas y circulares}
\label{subsec:matricial_ErEphi}

En el desarrollo de la teoría de difracción vectorial sobre dioptrías estigmáticas, es conveniente trabajar con las componentes cilíndricas del campo eléctrico $(E_r, E_\phi)$, pues estas se relacionan directamente con los modos de Fresnel a través de expresiones analíticas compactas (ver ecuación \eqref{eq:permisible_incident_field_given_Er_Ephi}).

Sin embargo, en la práctica experimental, las mediciones del campo eléctrico suelen realizarse en bases distintas: ya sea en componentes cartesianas $(E_x, E_y)$ mediante sensores planos orientados en coordenadas rectangulares, o en componentes circulares $(E_R, E_L)$ cuando se emplea polarimetría de luz polarizada circularmente.

En esta sección presentamos las transformaciones matriciales necesarias para expresar las componentes cilíndricas $(E_r, E_\phi)$ a partir de mediciones realizadas en estas bases alternativas.

\subsubsection{Transformación desde la base cartesiana}

La relación entre las componentes cartesianas $(E_x, E_y)$ y las componentes cilíndricas $(E_r, E_\phi)$ está dada por la matriz de rotación asociada al ángulo azimutal $\phi$:

\begin{equation}
\begin{bmatrix}
E_r \\[4pt]
E_\phi
\end{bmatrix}
=
\begin{bmatrix}
\cos\phi & \sin\phi \\[4pt]
-\sin\phi & \cos\phi
\end{bmatrix}
\begin{bmatrix}
E_x \\[4pt]
E_y
\end{bmatrix}.
\label{eq:Er_Ephi_from_cartesian}
\end{equation}

Esta transformación es particularmente útil cuando el campo se mide en un plano perpendicular al eje óptico mediante detectores alineados con los ejes $x$ e $y$.

\subsubsection{Transformación desde la base circular}

Las componentes de polarización circular se definen mediante los vectores de base:

\begin{equation}
\uvec{e}_R = \frac{1}{\sqrt{2}}(\uvec{x} - i\,\uvec{y}),
\qquad
\uvec{e}_L = \frac{1}{\sqrt{2}}(\uvec{x} + i\,\uvec{y}),
\label{eq:circular_basis_vectors}
\end{equation}

donde $\uvec{e}_R$ corresponde a la polarización circular derecha (right-handed) y $\uvec{e}_L$ a la polarización circular izquierda (left-handed).

La transformación desde componentes circulares $(E_R, E_L)$ a componentes cilíndricas $(E_r, E_\phi)$ está dada por:

\begin{equation}
\begin{bmatrix}
E_r \\[6pt]
E_\phi
\end{bmatrix}
=
\frac{1}{\sqrt{2}}
\begin{bmatrix}
e^{-i\phi} & e^{i\phi} \\[6pt]
-i\,e^{-i\phi} & i\,e^{i\phi}
\end{bmatrix}
\begin{bmatrix}
E_R \\[6pt]
E_L
\end{bmatrix}.
\label{eq:Er_Ephi_from_circular}
\end{equation}

Esta relación resulta especialmente relevante en el análisis de estados de polarización mediante la esfera de Poincaré y en experimentos donde se emplean láminas retardadoras de cuarto de onda para generar o analizar luz circularmente polarizada.

\subsubsection{Transformación inversa: de cilíndricas a cartesianas}

Por completitud, presentamos también la transformación inversa, que permite obtener las componentes cartesianas a partir de las cilíndricas:

\begin{equation}
\begin{bmatrix}
E_x \\[4pt]
E_y
\end{bmatrix}
=
\begin{bmatrix}
\cos\phi & -\sin\phi \\[4pt]
\sin\phi & \cos\phi
\end{bmatrix}
\begin{bmatrix}
E_r \\[4pt]
E_\phi
\end{bmatrix}.
\label{eq:cartesian_from_Er_Ephi}
\end{equation}

Estas transformaciones, en conjunto con la ecuación \eqref{eq:permisible_incident_field_given_Er_Ephi}, permiten determinar el campo incidente admisible en la base de Fresnel independientemente de la representación experimental empleada para su medición.
